% !TeX root = ../papers/nc-p3.tex
\section{種数5の有理水準における全境界円板}
\label{sec:rank-fifteen-pencil-trace-disks}
\begin{theorem}\label{thm:rank-fifteen-pencil-trace-disks}
$c=3/4$、$s=25/12$、$\xi=9/16$ の種数5水準の両表示を取る。
101で曲線の56境界点のいずれかに還元される標数零の通常窓は、
標準Koszul左右AS正則で、各始点で四中心の多項式環上の階数256の自由加群となる。
\end{theorem}
\begin{proof}
$K=k=\pm1$ 上では $u=n$ を使い、$D=\mathbb Z_{(101)}$、
$L^2=1+c^2$ とする。$K=k\sqrt{1+L^2u^2}$ とし、
$S^2=L^2+2cK+c^2u^2$、$T^2=L^2-2cK+c^2u^2$ について
$S(0)=c+k$、$T(0)=c-k$ の級数を取る。これらの初期値は単元である。
前節の関係から $r=K+cu^2+\sigma uS$、$w=K-cu^2+\tau uT$ である。
$K_0=\pm1$、$(\sigma,\tau)\in\{1,-1\}^2$ で全八点を尽くし、
二位置移動は $u$ の符号と二つの分岐符号を反転する。

\path{computations/verify_rank_fifteen_pencil_trace_disks.py} は、
二次関係の各Pl\"ucker座標を $D[[u]]/(u^{32})$ で評価する。
定数項に単元がない組は数体上での零を確認して $u$ で一度割り、
単元座標で行基底を取る。次に\emph{中心切断の前}に長さ5までの道を左に延ばす。
各位置で長さ1,2,3,4,5の文字次元を順に
$(1,1,1,1),(2,2,3,3),(5,5,5,5),(9,10,8,8),(14,14,14,14)$ として
付値環上で掃き出す。必要な階数に達するまで、単元ピボットがない段でだけ、
残る全行の定数項が有理数上で零であることを確認して $u$ で割る。
得られた語は通常持上げで全経路を張り、その総数は既知の長さ4,5の次元に等しいので
各文字の実際の基底を与える。

この基底で文字0,1の交換子方程式を作る。同様の飽和によって自由変数が
それぞれ3,2個になるまで消去する。これらは正の固定族の一般中心からの下界でもあるので、
残る方程式は分数体上、選んだ行の線型結合となる。単元ピボットを使う
後退代入は局所環の級数で行え、各通常持上げで同じ階数と核の自由変数を与える。
文字0で $(1,3,5),(2,7,11)$、文字1で二本の基底を選ぶと、
この手順で実際に持ち上がる四中心式が得られる。四中心商を右に延ばし、
定理\ref{thm:rank-fifteen-octic-all}と同じ上界を長さ13まで検証する。
$u$ で割った回数の最大値を各段の精度から引く。全32組で最後にも
正の精度が残るため、どの零判定にも十分な桁数がある。

残る有限境界では $K=0$ と $K=\pm s\pm1$ はこの同じ101で既に認証されている。
$K^2+sK-1=0$ では $r=\pm i$、$w=2K\pm\sqrt{4K^2-1}$、
$H=\pm\sqrt{K^2-1}$ を使う。上記コードは $\mathbb F_{101^4}$ に両根 $K$ と
三つの平方根の全選択を実現し、両表示で全四始点の四中心商を確認する。
最後に無限遠の八点では $L=\pm5/4$ に対応する $c+L=2,-1/2$ の両方で
定理\ref{thm:rank-fifteen-pencil-infinity}の重み付き関係を101で検算する。
六つの分岐値と他の有限境界値はこの標数でも所定の相異なる値を持つ。
以上と全次数の下界・中心商の上界から主張が従う。
\end{proof}
